\section{环境搭建与项目规范}

开始新的数据分析项目时，需要先设置好三个方面：

\begin{itemize}[nosep]
    \item \textbf{Environment}（环境）：Python 环境、包管理器、目录结构、输出约定
    \item \textbf{Skills}（技能）：把重复工作流（如 notebook 清理、spreadsheet 导出、最终报告打包）封装为可复用的 skill
    \item \textbf{Worktrees}（工作树）：将不同探索方向隔离到不同 worktree，避免假设、合并策略或可视化分支之间互相干扰
\end{itemize}

\subsection{用 AGENTS.md 引导 Codex 行为}

在接触数据之前，先告诉 Codex 项目的行为规范。个人默认设置放在 \texttt{\textasciitilde/.codex/AGENTS.md}，项目级规则放在仓库根目录的 \texttt{AGENTS.md}。

\begin{lstlisting}[caption={一个简洁的 AGENTS.md 示例},language={}]
## Data analysis defaults
- Use `uv run` or the project's existing Python environment.
- Keep source data in `data/raw/` and write cleaned data
  to `data/processed/`.
- Put exploratory notebooks in `analysis/` and final
  artifacts in `output/`.
- Never overwrite raw files.
- Prefer scripts or checked-in notebooks over unnamed
  scratch cells.
- Before merging datasets, report candidate keys, null
  rates, and join coverage.
\end{lstlisting}

\begin{importantbox}{AGENTS.md 的作用}
\texttt{AGENTS.md} 是 Codex 的``项目宪法''。它确保 Codex 在每次交互中都遵守统一的目录结构、数据保护规则和工作流约定，而不需要你在每条 prompt 里重复说明。
\end{importantbox}

\subsection{本章小结}
项目规范先行：通过 \texttt{AGENTS.md} 定义环境、目录结构和行为约束，通过 Skills 封装重复流程，通过 Worktrees 隔离并行探索。

